\scanpage{110}

\mspage{116}

\uline{Définition 2} Sous les conditions équivalentes de la \deleted{définition} prop. 1, on dit que $f$ est une $\mathbb D$\uline{-fibration cohomologique faible}. Sous les conditions duales, que $f$ est une $\mathbb D$\uline{-fibration homologique faible}.

C'est ce que dans la définition 1 p. 77 j'appelais ``$\mathbb D$-fibration cohomologique'' ou ``$\mathbb D$-fibration à gauche'', \inserted{et dualement.} Maintenant la notion me paraît trop faible pour mériter ce nom, que je réserve pour le cas ``complet''.

\uline{Prop 2} Les $\mathbb D$-fibrations coh. faibles sont stables par \unsure{changement|changements} de base lisses, les $\mathbb D$-fibr. hom. faibles stables par chang. de base propres.

Je devrais construire un exemple où \deleted{\unsure{on ait une}}\inserted{$f$ est une} $\mathbb D$-fibr. coh. faible, et ne l'est plus après chgt. de base.

\inserted{\uline{Prop. 3} \deleted{\unsure{Soit} $f:X\to Y$ de $\mathrm{Diag}$,} soit
\diagram{\begin{tikzcd}[ampersand replacement=\&,row sep=small,column sep=small]
 X\arrow[rr,"f"]\arrow[dr,"p"'] \& \& X'\arrow[dl,"p'"]\\
 \& S \&
\end{tikzcd}}
de $\mathrm{Diag}$, avec $p$ et $p'$ des $\mathbb D$-fibr. coh. faibles (resp. des $\mathbb D$-fibr. hom. faibles), soit $S_0\subset S$ tel que $\pi_0(S_0)\to\pi_0(S)$ surj. Pour que $f$ soit une $\mathbb D$-équiv., il suffit que \unsure{les} $f_s:X_s\to X'_s$, pour $s\in S_0$, le soient.}

\uline{Prop. 4} Si $f:X\to Y$ de $\mathrm{Diag}$ est $\mathbb D$-propre (resp. $\mathbb D$-lisse), alors $f$ est une $\mathbb D$-fibration coh. \deleted{faible} (resp. hom.) faible ssi pour tout $\unsure{F|E}$ constant sur $Y$, $f_* f^*F$ \inserted{(resp. $f_!f^*F$)} est loc. constant sur $Y$. Et s'il en est ainsi, \deleted{\unsure{D} cette condition est \unsure{co}} pour tout $Y'\to Y$ de $\mathrm{Diag}$, $f':X'=X\times_Y Y'\to Y'$ est encore une $\mathbb D$-fib. coh. (resp. hom.) faible.

\deleted{\unsure{En effet,}} Immédiat sur le critère b) de prop. 1.

% archive sheet number: 110
